%% Drafted from the mapped corpus by scripts/draft_topics.py.
% Model: gpt-5.6-terra; source characters: 15413.
\begin{konspekt}
\kreq{Дефиниция}{локален екстремум на функция на една променлива --- \kref{32.1}.}
\kreq{Формулировка и доказателство}{необходимо условие за локален екстремум при диференцируеми функции --- теорема на Ферма --- \kref{32.1}.}
\kreq{Доказателства}{теоремите на Рол, Лагранж и Коши, заедно с $g(a)\ne g(b)$ при предположенията на Коши --- \kref{32.2}.}
\kreq{Извеждане}{формула на Тейлър с остатъчен член във формата на Лагранж --- \kref{32.3}.}
\kextra{теоремата на Вайерщрас се \emph{използва без доказателство} при теоремата на Рол --- това е изрично разрешено от анотацията.}
\kreq{Литература}{[6], [8], [10].}
\end{konspekt}

\section{Локални екстремуми и теорема на Ферма}\label{s:32.1}

\begin{defn}[Локален екстремум]
Нека $f:D\to\mathbb{R}$, където $D\subseteq\mathbb{R}$, и нека $x_0\in D$.

Казваме, че $f$ има \emph{локален минимум} в точката $x_0$, ако съществува $\delta>0$, за което
\[
(x_0-\delta,x_0+\delta)\subseteq D
\]
и
\[
f(x_0)\leq f(x)
\quad\text{за всяко }x\in(x_0-\delta,x_0+\delta).
\]
Аналогично, $f$ има \emph{локален максимум} в $x_0$, ако съществува $\delta>0$, за което
\[
(x_0-\delta,x_0+\delta)\subseteq D
\]
и
\[
f(x_0)\geq f(x)
\quad\text{за всяко }x\in(x_0-\delta,x_0+\delta).
\]
Точка, в която функцията има локален минимум или локален максимум, се нарича \emph{точка на локален екстремум}.
\end{defn}

\begin{keythm}[Теорема на Ферма --- необходимо условие за локален екстремум]
Нека $f:D\to\mathbb{R}$ и нека $x_0$ е точка на локален екстремум на $f$. Ако $f$ е диференцируема в $x_0$, то
\[
f'(x_0)=0.
\]
\end{keythm}

\begin{proof}
Нека първо $x_0$ е точка на локален минимум. Тогава съществува $\delta>0$, такова че
\[
f(x_0)\leq f(x)
\]
за всяко $x\in(x_0-\delta,x_0+\delta)$.

Ако $x\in(x_0,x_0+\delta)$, то $x-x_0>0$ и $f(x)-f(x_0)\geq 0$. Следователно
\[
\frac{f(x)-f(x_0)}{x-x_0}\geq 0.
\]
При преминаване към граница отдясно получаваме
\[
f'(x_0)\geq 0.
\]

Ако $x\in(x_0-\delta,x_0)$, то $x-x_0<0$, докато отново $f(x)-f(x_0)\geq 0$. Затова
\[
\frac{f(x)-f(x_0)}{x-x_0}\leq 0.
\]
При преминаване към граница отляво следва
\[
f'(x_0)\leq 0.
\]
Тъй като производната съществува, левият и десният предел съвпадат. Следователно едновременно е изпълнено
\[
f'(x_0)\geq 0
\qquad\text{и}\qquad
f'(x_0)\leq 0,
\]
тоест $f'(x_0)=0$.

Ако $x_0$ е точка на локален максимум на $f$, то тя е точка на локален минимум на $-f$. От вече доказаното
\[
(-f)'(x_0)=0,
\]
откъдето $-f'(x_0)=0$, т.е. $f'(x_0)=0$.
\end{proof}

\begin{remark}
Теоремата на Ферма дава необходимо, но не достатъчно условие за локален екстремум. Условието $f'(x_0)=0$ само посочва възможна точка на екстремум.
\end{remark}

\section{Теореми за средните стойности}\label{s:32.2}

В този раздел нека $a<b$ и $f$ е непрекъсната в затворения интервал $[a,b]$ и диференцируема в отворения интервал $(a,b)$.

За доказателството на първата теорема използваме теоремата на Вайерщрас: всяка непрекъсната функция върху краен затворен интервал достига най-голямата и най-малката си стойност.

\subsection{Теорема на Рол}\label{s:32.2.1}

\begin{keythm}[Теорема на Рол]
Ако
\[
f(a)=f(b),
\]
то съществува точка $c\in(a,b)$, за която
\[
f'(c)=0.
\]
\end{keythm}

\begin{proof}
Понеже $f$ е непрекъсната в $[a,b]$, от теоремата на Вайерщрас следва, че тя достига своя минимум и максимум в някои точки
\[
x_{\min},x_{\max}\in[a,b].
\]

Ако поне една от точките $x_{\min}$ и $x_{\max}$ принадлежи на $(a,b)$, то в нея функцията има локален екстремум. Функцията е диференцируема в тази вътрешна точка, следователно по теоремата на Ферма производната в нея е нула. Получаваме търсената точка $c$.

Остава случаят, когато и минимумът, и максимумът се достигат само в краищата на интервала. Тъй като
\[
f(a)=f(b),
\]
минималната и максималната стойност са равни. Следователно функцията е константна в $[a,b]$. Тогава
\[
f'(x)=0
\quad\text{за всяко }x\in(a,b),
\]
и всяка точка от $(a,b)$ може да бъде избрана за $c$.
\end{proof}

\begin{supp}[Бележка]
В едно от предоставените изложения при разглеждането на вътрешна точка на минимум е записано $f(x_{\min})=0$. От теоремата на Ферма следва не равенство за стойността на функцията, а
\[
f'(x_{\min})=0.
\]
Именно това равенство се използва в доказателството на теоремата на Рол.
\end{supp}

\begin{remark}
Геометрично теоремата на Рол гласи, че ако графиката на функцията започва и завършва на една и съща височина, то между тези две точки има точка с хоризонтална допирателна.
\end{remark}

\subsection{Теорема на Лагранж}\label{s:32.2.2}

\begin{keythm}[Теорема на Лагранж за крайните нараствания]
Съществува точка $c\in(a,b)$, такава че
\[
f(b)-f(a)=f'(c)(b-a).
\]
Еквивалентно,
\[
f'(c)=\frac{f(b)-f(a)}{b-a}.
\]
\end{keythm}

\begin{proof}
Да разгледаме функцията
\[
h(x)=f(x)-kx,
\]
където константата $k$ ще изберем така, че да е изпълнено условието на теоремата на Рол:
\[
h(a)=h(b).
\]
Това равенство е еквивалентно на
\[
f(a)-ka=f(b)-kb,
\]
или, тъй като $a\ne b$,
\[
k=\frac{f(b)-f(a)}{b-a}.
\]

Функцията $h$ е непрекъсната в $[a,b]$ и диференцируема в $(a,b)$, защото $f$ има същите свойства. При избраното $k$ е изпълнено $h(a)=h(b)$. По теоремата на Рол съществува $c\in(a,b)$, за което
\[
h'(c)=0.
\]
Но
\[
h'(x)=f'(x)-k.
\]
Следователно
\[
0=h'(c)=f'(c)-k,
\]
тоест
\[
f'(c)=k=\frac{f(b)-f(a)}{b-a}.
\]
След умножение по $b-a$ получаваме
\[
f(b)-f(a)=f'(c)(b-a).
\]
\end{proof}

\begin{remark}
Числото
\[
\frac{f(b)-f(a)}{b-a}
\]
е средната скорост на изменение на функцията върху $[a,b]$, а $f'(c)$ е моментната скорост на изменение в подходяща вътрешна точка $c$.
\end{remark}

\subsection{Теорема на Коши}\label{s:32.2.3}

\begin{keythm}[Обобщена теорема на Лагранж --- теорема на Коши]
Нека $f$ и $g$ са непрекъснати в $[a,b]$ и диференцируеми в $(a,b)$. Нека освен това
\[
g'(x)\ne 0
\quad\text{за всяко }x\in(a,b).
\]
Тогава съществува $c\in(a,b)$, такова че
\[
\frac{f(b)-f(a)}{g(b)-g(a)}
=
\frac{f'(c)}{g'(c)}.
\]
\end{keythm}

\begin{proof}
\emph{Стъпка 1: знаменателят е ненулев.} Да допуснем противното, тоест $g(a)=g(b)$. Тогава $g$ е непрекъсната в $[a,b]$, диференцируема в $(a,b)$ и приема равни стойности в двата края. По теоремата на Рол съществува $\xi\in(a,b)$, за което
\[
g'(\xi)=0.
\]
Това противоречи на условието $g'(x)\ne 0$ за всяко $x\in(a,b)$. Следователно
\[
g(b)-g(a)\ne 0,
\]
така че лявата страна на доказваното равенство е определена.

\emph{Стъпка 2: спомагателната функция.} Полагаме
\[
H(x)=\bigl(f(b)-f(a)\bigr)\bigl(g(x)-g(a)\bigr)-\bigl(g(b)-g(a)\bigr)\bigl(f(x)-f(a)\bigr).
\]
Числата $f(b)-f(a)$ и $g(b)-g(a)$ са константи, затова $H$ е линейна комбинация на $f$ и $g$ с добавена константа. Значи $H$ е непрекъсната в $[a,b]$ и диференцируема в $(a,b)$, защото $f$ и $g$ имат тези свойства.

\emph{Стъпка 3: равни стойности в краищата.} Пресмятаме двете стойности:
\[
H(a)=\bigl(f(b)-f(a)\bigr)\underbrace{\bigl(g(a)-g(a)\bigr)}_{=0}-\bigl(g(b)-g(a)\bigr)\underbrace{\bigl(f(a)-f(a)\bigr)}_{=0}=0,
\]
\[
H(b)=\bigl(f(b)-f(a)\bigr)\bigl(g(b)-g(a)\bigr)-\bigl(g(b)-g(a)\bigr)\bigl(f(b)-f(a)\bigr)=0,
\]
тъй като двете произведения в $H(b)$ съвпадат. Следователно
\[
H(a)=H(b)=0.
\]

\emph{Стъпка 4: прилагане на теоремата на Рол.} Функцията $H$ удовлетворява трите условия на Рол, затова съществува $c\in(a,b)$, за което
\[
H'(c)=0.
\]

\emph{Стъпка 5: разписване на производната.} Диференцираме $H$, като константните множители излизат пред производната, а производните на $g(x)-g(a)$ и $f(x)-f(a)$ са съответно $g'(x)$ и $f'(x)$:
\[
H'(x)=\bigl(f(b)-f(a)\bigr)g'(x)-\bigl(g(b)-g(a)\bigr)f'(x).
\]
Полагаме $x=c$ и използваме $H'(c)=0$:
\[
0=\bigl(f(b)-f(a)\bigr)g'(c)-\bigl(g(b)-g(a)\bigr)f'(c),
\]
тоест
\[
\bigl(f(b)-f(a)\bigr)g'(c)=\bigl(g(b)-g(a)\bigr)f'(c).
\]

\emph{Стъпка 6: деление.} По условие $g'(c)\ne 0$, а по стъпка 1 $g(b)-g(a)\ne 0$. Следователно произведението
\[
\bigl(g(b)-g(a)\bigr)g'(c)\ne 0
\]
и можем да разделим на него двете страни на последното равенство:
\[
\frac{\bigl(f(b)-f(a)\bigr)g'(c)}{\bigl(g(b)-g(a)\bigr)g'(c)}
=
\frac{\bigl(g(b)-g(a)\bigr)f'(c)}{\bigl(g(b)-g(a)\bigr)g'(c)}.
\]
Отляво се съкращава $g'(c)$, а отдясно --- $g(b)-g(a)$, откъдето
\[
\frac{f(b)-f(a)}{g(b)-g(a)}
=
\frac{f'(c)}{g'(c)}.
\]
\end{proof}

\begin{prop}
При условията на теоремата на Коши е изпълнено
\[
g(a)\ne g(b).
\]
\end{prop}

\begin{proof}
Да допуснем противното:
\[
g(a)=g(b).
\]
Тогава функцията $g$ удовлетворява условията на теоремата на Рол. Следователно съществува $c\in(a,b)$, за което
\[
g'(c)=0.
\]
Това противоречи на предположението, че $g'(x)\ne 0$ за всяко $x\in(a,b)$. Следователно
\[
g(a)\ne g(b).
\]
\end{proof}

\begin{proof}[Второ доказателство на теоремата на Коши]
От предходното твърдение знаменателят $g(b)-g(a)$ е различен от нула. Нека
\[
k=\frac{f(b)-f(a)}{g(b)-g(a)}
\]
и дефинираме
\[
h(x)=f(x)-k g(x).
\]
Като линейна комбинация на $f$ и $g$, функцията $h$ е непрекъсната в $[a,b]$ и диференцируема в $(a,b)$.

От избора на $k$ следва
\[
k\bigl(g(b)-g(a)\bigr)=f(b)-f(a).
\]
Следователно
\[
f(a)-kg(a)=f(b)-kg(b),
\]
тоест
\[
h(a)=h(b).
\]
По теоремата на Рол съществува $c\in(a,b)$, за което
\[
h'(c)=0.
\]
Тъй като
\[
h'(x)=f'(x)-k g'(x),
\]
получаваме
\[
f'(c)-k g'(c)=0.
\]
Понеже $g'(c)\ne 0$, можем да разделим на $g'(c)$:
\[
\frac{f'(c)}{g'(c)}=k.
\]
От определението на $k$ следва
\[
\frac{f'(c)}{g'(c)}
=
\frac{f(b)-f(a)}{g(b)-g(a)}.
\]
\end{proof}

\begin{cor}
Теоремата на Лагранж е частен случай на теоремата на Коши при
\[
g(x)=x.
\]
\end{cor}

\begin{proof}
За $g(x)=x$ имаме $g'(x)=1$ и $g(b)-g(a)=b-a$. Формулата на Коши приема вида
\[
\frac{f(b)-f(a)}{b-a}=f'(c).
\]
\end{proof}

\section{Формула на Тейлър}\label{s:32.3}

\subsection{Полином на Тейлър}\label{s:32.3.1}

\begin{defn}[Полином на Тейлър]
Нека функцията $f$ притежава производни до ред $n$ в точката $x_0$. Полиномът
\[
T_n(x)=\sum_{k=0}^{n}\frac{f^{(k)}(x_0)}{k!}(x-x_0)^k
\]
се нарича \emph{полином на Тейлър от ред $n$} на функцията $f$ около точката $x_0$.

Разписан, той има вида
\[
T_n(x)=f(x_0)+f'(x_0)(x-x_0)
+\frac{f''(x_0)}{2!}(x-x_0)^2+\cdots
+\frac{f^{(n)}(x_0)}{n!}(x-x_0)^n.
\]
\end{defn}

\begin{lem}
Ако $f$ притежава производни до ред $n$ в $x_0$, то
\[
T_n^{(i)}(x_0)=f^{(i)}(x_0),
\qquad i=0,1,\ldots,n.
\]
\end{lem}

\begin{proof}
При диференциране $i$ пъти всеки член
\[
\frac{f^{(k)}(x_0)}{k!}(x-x_0)^k
\]
с $k<i$ изчезва. За $k\geq i$ се получава множител
\[
k(k-1)\cdots(k-i+1)(x-x_0)^{k-i}.
\]
При $x=x_0$ всички членове с $k>i$ са нула, защото съдържат положителна степен на $x-x_0$. Остава само членът с $k=i$, който е равен на
\[
f^{(i)}(x_0).
\]
Следователно
\[
T_n^{(i)}(x_0)=f^{(i)}(x_0).
\]
\end{proof}

\subsection{Остатъчен член във формата на Лагранж}\label{s:32.3.2}

Полиномът на Тейлър приближава $f$ около $x_0$; разликата между функцията и приближението получава собствено име.

\begin{defn}[Остатъчен член]
Нека $f$ притежава производни до ред $n$ в $x_0$ и нека $T_n$ е нейният полином на Тейлър около $x_0$. Функцията
\[
R_n(x)=f(x)-T_n(x)
\]
се нарича \emph{остатъчен член от ред $n$}. По построение
\[
f(x)=T_n(x)+R_n(x).
\]
\end{defn}

От лемата по-горе $T_n^{(i)}(x_0)=f^{(i)}(x_0)$ за $i=0,1,\ldots,n$, следователно
\[
R_n^{(i)}(x_0)=0,\qquad i=0,1,\ldots,n.
\]
Остатъкът и първите му $n$ производни се анулират в $x_0$; въпросът е с каква скорост той расте при отдалечаване от $x_0$. Теоремата по-долу дава точен отговор чрез стойност на $(n+1)$-вата производна във вътрешна точка.

\begin{keythm}[Формула на Тейлър с остатъчен член на Лагранж]
Нека са дадени цяло число $n\geq 0$, точка $x_0\in\mathbb{R}$ и число $\delta>0$. Нека функцията
\[
f:(x_0-\delta,x_0+\delta)\to\mathbb{R}
\]
притежава производни до ред $n+1$ включително във \emph{всяка} точка на интервала, тоест $f^{(n+1)}(t)$ съществува за всяко $t\in(x_0-\delta,x_0+\delta)$.

Тогава за всяко
\[
x\in(x_0-\delta,x_0+\delta),\qquad x\ne x_0,
\]
съществува точка $c$, разположена \emph{строго} между $x_0$ и $x$, тоест
\[
c\in(x_0,x)\ \text{ при } x>x_0,
\qquad
c\in(x,x_0)\ \text{ при } x<x_0,
\]
за която
\[
f(x)
=
\sum_{k=0}^{n}\frac{f^{(k)}(x_0)}{k!}(x-x_0)^k
+
\frac{f^{(n+1)}(c)}{(n+1)!}(x-x_0)^{n+1}.
\]
Еквивалентно, съществува $\theta\in(0,1)$, за което $c=x_0+\theta(x-x_0)$.

При $x=x_0$ равенството е изпълнено тривиално: сумата се свежда до $f(x_0)$, а остатъчният член е нула.
\end{keythm}

Условието е върху отворен интервал около $x_0$, а не върху затворен: така за всяко $x$ от интервала целият затворен интервал с краища $x_0$ и $x$ лежи вътре в областта, върху която $f^{(n+1)}$ съществува. От съществуването на $f^{(k+1)}$ следва непрекъснатост на $f^{(k)}$, затова $f,f',\ldots,f^{(n)}$ са непрекъснати в целия интервал и отделно предположение за това не е нужно.

\begin{proof}
Случаят $x=x_0$ е непосредствен: всички степени $(x-x_0)^k$ при $k\geq 1$ са нула, така че и двете страни са равни на $f(x_0)$.

Нека по-нататък $x\ne x_0$ е фиксирано. Означаваме с
\[
I=[\min(x_0,x),\ \max(x_0,x)]
\]
затворения интервал с краища $x_0$ и $x$. Тъй като и двете точки лежат в $(x_0-\delta,x_0+\delta)$, а този интервал е изпъкнал, имаме $I\subseteq(x_0-\delta,x_0+\delta)$. Достатъчно е да докажем, че
\[
R_n(x)=\frac{f^{(n+1)}(c)}{(n+1)!}(x-x_0)^{n+1}
\]
за подходяща вътрешна точка $c$ на $I$, защото $f(x)=T_n(x)+R_n(x)$.

\emph{Стъпка 1: спомагателните функции.} При фиксирано $x$ разглеждаме функциите на променливата $t$
\[
\varphi(t)
=
f(x)-\sum_{k=0}^{n}\frac{f^{(k)}(t)}{k!}(x-t)^k
=
f(x)-f(t)-\frac{f'(t)}{1!}(x-t)-\cdots-\frac{f^{(n)}(t)}{n!}(x-t)^n
\]
и
\[
\psi(t)=(x-t)^{n+1}.
\]
Обърнете внимание, че тук $x$ е константа, а $t$ пробягва $I$; в $\varphi$ производните на $f$ се пресмятат в $t$, а не в $x_0$.

Функцията $\varphi$ е крайна сума от произведения на $f^{(k)}(t)$ с полиноми на $t$. За $k\leq n$ производната $f^{(k)}$ е диференцируема в целия интервал $(x_0-\delta,x_0+\delta)$, защото $f^{(k+1)}$ съществува там. Следователно $\varphi$ е диференцируема, а значи и непрекъсната, в отворено множество, съдържащо $I$. Същото важи очевидно и за полинома $\psi$. В частност $\varphi$ и $\psi$ са непрекъснати в $I$ и диференцируеми във вътрешността му.

\emph{Стъпка 2: стойности в краищата.} При $t=x$ всички степени $(x-t)^k$ с $k\geq 1$ се анулират, така че от сумата остава само членът с $k=0$:
\[
\varphi(x)=f(x)-f(x)=0.
\]
При $t=x_0$ сумата е точно полиномът на Тейлър:
\[
\varphi(x_0)
=
f(x)-\sum_{k=0}^{n}\frac{f^{(k)}(x_0)}{k!}(x-x_0)^k
=
f(x)-T_n(x)
=
R_n(x).
\]
Аналогично
\[
\psi(x)=0,
\qquad
\psi(x_0)=(x-x_0)^{n+1}\ne 0.
\]

\emph{Стъпка 3: разписване на $\varphi'$.} Диференцираме почленно, като за $k\geq 1$ използваме правилото за произведение:
\[
\frac{d}{dt}\left[\frac{f^{(k)}(t)}{k!}(x-t)^k\right]
=
\frac{f^{(k+1)}(t)}{k!}(x-t)^k
-
\frac{f^{(k)}(t)}{k!}\,k\,(x-t)^{k-1},
\]
където вторият член идва от $\frac{d}{dt}(x-t)^k=-k(x-t)^{k-1}$. За $k=0$ вторият член липсва, тъй като $(x-t)^0=1$ е константа. Понеже $\dfrac{k}{k!}=\dfrac{1}{(k-1)!}$, получаваме
\[
\varphi'(t)
=
-\sum_{k=0}^{n}\frac{f^{(k+1)}(t)}{k!}(x-t)^k
+
\sum_{k=1}^{n}\frac{f^{(k)}(t)}{(k-1)!}(x-t)^{k-1}.
\]
Във втората сума полагаме $j=k-1$ и тя се превръща в
\[
\sum_{j=0}^{n-1}\frac{f^{(j+1)}(t)}{j!}(x-t)^{j},
\]
тоест в първата сума без нейния последен член. Двете суми се унищожават член по член и остава
\[
\varphi'(t)
=
-\frac{f^{(n+1)}(t)}{n!}(x-t)^n.
\]
Това е причината да се избере точно тази спомагателна функция: сумата е телескопична и нейната производна съдържа само $(n+1)$-вата производна на $f$.

За $\psi$ имаме направо
\[
\psi'(t)=-(n+1)(x-t)^n.
\]
За всяко $t$ от вътрешността на $I$ е $t\ne x$, следователно $(x-t)^n\ne 0$ и
\[
\psi'(t)\ne 0.
\]

\emph{Стъпка 4: прилагане на теоремата на Коши.} Функциите $\varphi$ и $\psi$ са непрекъснати в $I$, диференцируеми във вътрешността му и $\psi'$ не се анулира там. По теоремата на Коши съществува вътрешна точка $c$ на $I$, тоест точка строго между $x_0$ и $x$, за която
\[
\frac{\varphi(x)-\varphi(x_0)}
{\psi(x)-\psi(x_0)}
=
\frac{\varphi'(c)}{\psi'(c)}.
\]
Теоремата на Коши е формулирана за интервал $[a,b]$ с $a<b$. Ако $x<x_0$, тя се прилага с $a=x$ и $b=x_0$; лявата страна не се променя, защото при размяна на краищата числителят и знаменателят сменят знака си едновременно.

\emph{Стъпка 5: заместване.} Използваме стойностите от стъпка 2 и производните от стъпка 3:
\[
\frac{0-R_n(x)}{0-(x-x_0)^{n+1}}
=
\frac{
-\dfrac{f^{(n+1)}(c)}{n!}(x-c)^n
}{
-(n+1)(x-c)^n
}.
\]
Отляво знаците се съкращават. Отдясно $c$ е вътрешна точка, затова $c\ne x$ и $(x-c)^n\ne 0$, така че този множител може да бъде съкратен; освен това
\[
\frac{1}{n!\,(n+1)}=\frac{1}{(n+1)!}.
\]
Получаваме
\[
\frac{R_n(x)}{(x-x_0)^{n+1}}
=
\frac{f^{(n+1)}(c)}{(n+1)!},
\]
а след умножение по $(x-x_0)^{n+1}\ne 0$
\[
R_n(x)
=
\frac{f^{(n+1)}(c)}{(n+1)!}(x-x_0)^{n+1}.
\]
Тъй като $f(x)=T_n(x)+R_n(x)$, това е точно формулата на Тейлър с остатъчен член на Лагранж.
\end{proof}

\begin{remark}
Точката $c$ в остатъчния член зависи както от $x$, така и от $n$; тя не е дадена явно, а само съществува. Затова често се означава с $\xi(x)$:
\[
f(x)
=
f(x_0)+f'(x_0)(x-x_0)+\cdots+
\frac{f^{(n)}(x_0)}{n!}(x-x_0)^n
+
\frac{f^{(n+1)}(\xi(x))}{(n+1)!}(x-x_0)^{n+1}.
\]
Винаги $\xi(x)$ е между $x_0$ и $x$.
\end{remark}

\begin{example}
За $n=1$ формулата на Тейлър приема вида
\[
f(x)=f(x_0)+f'(x_0)(x-x_0)
+\frac{f''(c)}{2}(x-x_0)^2,
\]
където $c$ е между $x_0$ и $x$.
\end{example}

\section{Изпитен фокус}\label{s:32.4}

\begin{itemize}
\item Да се дефинират локален минимум, локален максимум и локален екстремум, като се подчертае съществуването на околност на точката, съдържаща се в дефиниционната област.

\item Да се формулира и докаже теоремата на Ферма чрез знака на отношението
\[
\frac{f(x)-f(x_0)}{x-x_0}
\]
отляво и отдясно на точка на локален минимум; случаят на локален максимум се свежда до функцията $-f$.

\item Да се формулира теоремата на Рол и да се възпроизведе доказателството чрез теоремата на Вайерщрас и теоремата на Ферма. Трябва да се разгледа и случаят, когато минимумът и максимумът се достигат само в краищата.

\item Да се формулира теоремата на Лагранж и да се изведе от Рол чрез спомагателната функция
\[
h(x)=f(x)-kx,
\qquad
k=\frac{f(b)-f(a)}{b-a}.
\]

\item Да се формулира теоремата на Коши, включително условието
\[
g'(x)\ne 0\quad\text{за }x\in(a,b).
\]
Да се докаже предварително, че то влече $g(a)\ne g(b)$, чрез противоречие с теоремата на Рол.

\item Да се докаже теоремата на Коши чрез функцията
\[
h(x)=f(x)-k g(x),
\qquad
k=\frac{f(b)-f(a)}{g(b)-g(a)}.
\]

\item Да се знае полиномът на Тейлър
\[
T_n(x)=\sum_{k=0}^{n}\frac{f^{(k)}(x_0)}{k!}(x-x_0)^k
\]
и свойството, че неговите производни до ред $n$ в $x_0$ съвпадат със съответните производни на $f$.

\item Да се формулира формулата на Тейлър с остатъчен член във формата на Лагранж: предположението е за съществуване на $f^{(n+1)}$ във всяка точка на отворен интервал около $x_0$, а точката $c$ е строго между $x_0$ и $x$.

\item Да се възпроизведе доказателството: фиксира се $x$, въвеждат се
\[
\varphi(t)=f(x)-\sum_{k=0}^{n}\frac{f^{(k)}(t)}{k!}(x-t)^k,
\qquad
\psi(t)=(x-t)^{n+1},
\]
пресмятат се $\varphi(x)=0$, $\varphi(x_0)=R_n(x)$, $\psi(x)=0$, $\psi(x_0)=(x-x_0)^{n+1}$, разписва се телескопичната производна
\[
\varphi'(t)=-\frac{f^{(n+1)}(t)}{n!}(x-t)^n,
\]
и към $\varphi$ и $\psi$ се прилага теоремата на Коши.
\end{itemize}
